\begin{notechapter}
\chapter{Transfer Learning}\label{ch:transfer-learning}
\chaptersubtitle{Reusing pretrained networks: preprocessing, feature extraction and fine-tuning}

\begin{chapterintro}
Training a large CNN from scratch needs a large labelled data set. Most applications do
not have one, but they can start from a network pretrained on ImageNet: its features are
useful far beyond the classes it was trained on. This chapter explains why features
transfer, how to prepare inputs exactly as the pretrained network expects, and the two
ways of reusing it: as a fixed feature extractor, or as an initialization for fine-tuning.
You will learn
\begin{itemize}[nosep]
  \item how to standardize images per channel, and how to compute the statistics correctly;
  \item what changes from the first to the last layers of a CNN, and what this implies for
        transfer;
  \item feature extraction with a new linear head or a classical classifier, without data
        leakage;
  \item fine-tuning: which layers, which learning rates, and what to do with batch
        normalization;
  \item when transfer helps, and when it does not.
\end{itemize}
\prerequisites{\cref{ch:cnn-architectures}; optimizers and weight decay (\cref{ch:optimizers});
batch normalization (\cref{ch:normalization}); validation and data leakage
(\cref{ch:measuring-performance}).}
\end{chapterintro}

\section{Preprocessing images}\label{transfer-learning:sec:preprocessing}

\paragraph{Why standardize.} Pixel values are stored as integers in $[0, 255]$. Gradient
descent on the first layer converges faster when its inputs are centred and of unit scale:
inputs with a large common mean make the first layer badly conditioned
(\cref{activation-functions:sec:centred}), and their scale multiplies the effective learning
rate of the first-layer weights.

\paragraph{The options.}
\begin{itemize}
  \item scaling to $[0,1]$: $x/255$ (\texttt{transforms.ToTensor()} does it);
  \item scaling to $[-1,1]$: $(x - 127.5)/127.5$, common in generative models;
  \item \emph{per-channel standardization}: $x^{(c)}\mapsto(x^{(c)} - \mu_c)/\sigma_c$ for each channel $c$, where $\mu_c$ and $\sigma_c$ are the
        mean and standard deviation of channel $c$ over all pixels of all training images.
\end{itemize}
Per-channel statistics respect the different distributions of the colour channels; one
global mean and standard deviation would not. The ImageNet values on $[0,1]$-scaled
images, $\mu = (0.485, 0.456, 0.406)$ and $\sigma = (0.229, 0.224, 0.225)$ for (R, G, B), are those used by the
\texttt{torchvision} models.

\begin{pitfall}[Statistics from the training set, applied everywhere]
Compute $\mu_c$ and $\sigma_c$ on the training set only, and apply the same values to validation and
test data: recomputing them on test data, or on all the data, leaks information about the
evaluation set into the model (\cref{ch:measuring-performance}). With a \emph{pretrained}
network, use the statistics it was trained with (and its resizing and cropping), not those of
your data set: its first layer expects exactly that input distribution.
\end{pitfall}

\paragraph{Computing the statistics correctly.} The standard deviation of a channel over the
data set is \emph{not} the average of the per-image standard deviations. By the law of total
variance, the variance over all pixels is the average within-image variance \emph{plus} the
variance of the image means: averaging per-image standard deviations underestimates it.
Accumulate sums and sums of squares instead:

\begin{lstlisting}[language=Python]
import torch
from torch.utils.data import DataLoader

def channel_mean_std(dataset, batch_size=256):
    """Per-channel mean and std over all pixels of all images (pooled)."""
    loader = DataLoader(dataset, batch_size=batch_size, shuffle=False)
    n, s, s2 = 0, 0.0, 0.0
    for images, _ in loader:                    # images: (B, C, H, W) in [0, 1]
        x = images.transpose(0, 1).flatten(1)   # (C, B*H*W)
        n += x.shape[1]
        s = s + x.sum(dim=1)
        s2 = s2 + (x ** 2).sum(dim=1)
    mean = s / n
    std = (s2 / n - mean ** 2).sqrt()
    return mean, std
\end{lstlisting}

On $50$ random images of varying contrast, this gives $\sigma\approx0.220$ per channel, the exact value,
while the average of per-image standard deviations gives $0.171$. (For very large data sets in
low precision, accumulate in double precision or use Welford's algorithm.)

\begin{lstlisting}[language=Python]
from torchvision import datasets, transforms

mean, std = channel_mean_std(datasets.CIFAR10("data", train=True, download=True,
                                              transform=transforms.ToTensor()))
train_tf = transforms.Compose([transforms.RandomCrop(32, padding=4),       # augmentation:
                               transforms.RandomHorizontalFlip(),          # training only
                               transforms.ToTensor(), transforms.Normalize(mean, std)])
test_tf = transforms.Compose([transforms.ToTensor(), transforms.Normalize(mean, std)])
train_set = datasets.CIFAR10("data", train=True, transform=train_tf)
test_set = datasets.CIFAR10("data", train=False, transform=test_tf)
# a validation set is split off train_set (with test_tf), never taken from test_set
\end{lstlisting}

\paragraph{Data augmentation.} Random crops, flips, small rotations or colour changes
create new training images with the same label: they encode invariances the model should
have and act as a regularizer. They are applied to training images only; evaluation uses
the deterministic transform.

\section{Why features transfer}\label{transfer-learning:sec:why}

A CNN trained for classification computes a hierarchy of features. Visualizations show
oriented edges and colour blobs in the first layer, textures and simple patterns in the
next, parts of objects in the middle, and class-specific configurations near the top
(Zeiler and Fergus, \emph{Visualizing and Understanding Convolutional Networks}, ECCV,
2014). The first layers are \emph{general}: similar features are learned on any natural images.
The last layers are \emph{specific} to the training task. Yosinski, Clune, Bengio and Lipson (\emph{How
Transferable Are Features in Deep Neural Networks?}, NeurIPS, 2014) measured this by
transferring the first $n$ layers between networks trained on two halves of ImageNet:
transfer degrades as $n$ grows and as the tasks become more different, and initializing with
transferred features and fine-tuning helped even when the new data set was large.
The penultimate features of an ImageNet network are, in practice, strong generic image
descriptors: a linear classifier on them is a hard baseline for many vision tasks (Razavian,
Azizpour, Sullivan and Carlsson, \emph{CNN Features Off-the-Shelf}, CVPR Workshops, 2014).

\section{Feature extraction}\label{transfer-learning:sec:extraction}

Remove the network's last (task-specific) layer, freeze the rest, and use the output of global
average pooling as a fixed feature vector: $512$ numbers for ResNet-18, $2048$ for ResNet-50. Then
either train a new linear layer on top (\emph{linear probing}), or extract the features once and
train any classifier on them.

\begin{lstlisting}[language=Python]
import torch
from torch import nn
from torchvision import models

model = models.resnet18(weights=models.ResNet18_Weights.IMAGENET1K_V1)
model.eval()                                       # batch norm in inference mode
extractor = nn.Sequential(*list(model.children())[:-1], nn.Flatten())

@torch.no_grad()
def features(loader):
    feats, labels = [], []
    for x, y in loader:                            # x preprocessed as for ImageNet
        feats.append(extractor(x)); labels.append(y)
    return torch.cat(feats).numpy(), torch.cat(labels).numpy()

F_train, y_train = features(train_loader)          # (n_train, 512)
F_test, y_test = features(test_loader)
\end{lstlisting}

Any classifier can then be trained on the features: logistic regression, an SVM
(\cref{ch:svm}), a random forest (\cref{ch:ensemble-methods}). With a few hundred examples
and $512$ or $2048$ features, dimensionality reduction by PCA (\cref{ch:pca-overview}) can help.
Every fitted step, the standardization and the PCA included, must be fitted on training data
only; a pipeline guarantees it:

\begin{lstlisting}[language=Python]
from sklearn.pipeline import make_pipeline
from sklearn.preprocessing import StandardScaler
from sklearn.decomposition import PCA
from sklearn.svm import SVC

clf = make_pipeline(StandardScaler(), PCA(n_components=100), SVC(kernel="rbf", C=1.0))
clf.fit(F_train, y_train)                         # scaler, PCA and SVM fitted on train only
print(f"test accuracy: {clf.score(F_test, y_test):.3f}")
\end{lstlisting}

Fitting the PCA on all the features before splitting, a common mistake, uses the test set to
choose the projection. Hyperparameters ($C$, the number of components) are tuned by
cross-validation on the training set.

\section{Fine-tuning}\label{transfer-learning:sec:finetuning}

Fine-tuning continues training the pretrained weights on the new task, starting from the
pretrained values, together with the new head.

\begin{lstlisting}[language=Python]
model = models.resnet18(weights=models.ResNet18_Weights.IMAGENET1K_V1)
for p in model.parameters():                       # freeze everything ...
    p.requires_grad = False
model.fc = nn.Linear(model.fc.in_features, 5)      # ... except a new head for 5 classes
for p in model.layer4.parameters():                # ... and the last stage
    p.requires_grad = True
opt = torch.optim.AdamW([
    {"params": model.layer4.parameters(), "lr": 1e-4},
    {"params": model.fc.parameters(), "lr": 1e-3}], weight_decay=0.01)
\end{lstlisting}

\paragraph{Which layers.} The less data, and the closer the new task to ImageNet, the fewer
layers should be trained. Typical choices: only the head (linear probing) for very small data
sets; the last stage and the head for small ones; everything for large ones. Above, $8.4$ of the
$11.2$ million parameters are trainable.

\paragraph{Learning rates.} Pretrained layers need small steps, or the first updates, driven
by a randomly initialized head, destroy the features. Common practice is a learning rate
about ten times smaller than for training from scratch, smaller still for earlier layers
(\emph{discriminative} learning rates, as in the parameter groups above), with warm-up
(\cref{optimizers:sec:schedules}). Training the head first with the backbone frozen, then
fine-tuning everything (``linear probing then fine-tuning'') protects the features; Kumar,
Raghunathan, Jones, Ma and Liang (\emph{Fine-Tuning can Distort Pretrained Features and
Underperform Out-of-Distribution}, ICLR, 2022) show that it also improves robustness to
distribution shift.

\paragraph{Batch normalization.} The running statistics of the pretrained network describe
ImageNet activations. With reasonable batches they adapt during fine-tuning; with very
small batches, the batch statistics are noisy (\cref{normalization:sec:bn}), and it is better to
keep the batch normalization layers in evaluation mode (frozen statistics), or to freeze them
entirely.

\section{When does transfer help?}\label{transfer-learning:sec:when}

\begin{center}\small
\begin{tabular}{lp{4.6cm}p{4.6cm}}
\toprule
& similar task to the source & different task \\
\midrule
small data set & feature extraction, linear head & features of earlier layers; careful fine-tuning \\
large data set & fine-tune everything & fine-tune everything, or train from scratch \\
\bottomrule
\end{tabular}
\end{center}

Transfer is most useful with little labelled data. For images very different from natural
photographs (medical scans, satellite multispectral images, spectrograms), the general
early features still help, but the higher ones less, and \emph{negative transfer}, where
pretraining hurts, is possible. Inputs with a different number of channels need an adapted
first layer (e.g.\ averaging the RGB filters for grey-level images). The same pattern of
pretraining on a large generic data set and adapting to a small specific one is the basis of
language models, with BERT and its successors in place of ImageNet CNNs
(\cref{part:language}).

\begin{chaptersummary}
\begin{itemize}[nosep]
  \item Standardize inputs per channel with statistics computed on all pixels of the training
        set, or, for a pretrained network, with the statistics and resizing it was trained with.
        Averaging per-image standard deviations underestimates the true one.
  \item CNN features go from general (edges, textures) to task-specific; penultimate features of
        ImageNet networks are strong generic descriptors.
  \item Feature extraction freezes the backbone and trains a head or a classical classifier;
        every fitted preprocessing step must be fitted on training data only.
  \item Fine-tuning trains some or all pretrained layers with small, often layer-dependent,
        learning rates; probing first protects the features; batch normalization needs care
        with small batches.
  \item Transfer helps most with little data and similar domains; it can hurt for very different
        inputs.
\end{itemize}
\end{chaptersummary}

\section*{Exercises}
\addcontentsline{toc}{section}{Exercises}

\begin{exercise}\label{transfer-learning:exr:1}
Two grey-level images of $2$ pixels: $(0.1, 0.3)$ and $(0.6, 0.8)$. Compute the pooled standard deviation of the
pixels and the average of the two per-image standard deviations, and check the law of total
variance.
\solutionlink{transfer-learning:exr:1}
\end{exercise}

\begin{exercise}\label{transfer-learning:exr:2}
A first layer computes $z = w^\top x + b$. Show that standardizing a channel, $x\mapsto(x - \mu)/\sigma$, can be undone by
changing $w$ and $b$. Why does it still matter for training?
\solutionlink{transfer-learning:exr:2}
\end{exercise}

\begin{exercise}\label{transfer-learning:exr:3}
Count the parameters of ResNet-18's new head for $5$ classes, and those of its last stage
(\texttt{layer4}), given that ResNet-18 has $11\,689\,512$ parameters with its original $1000$-class head and $8\,396\,293$ are
trainable in the fine-tuning example.
\solutionlink{transfer-learning:exr:3}
\end{exercise}

\begin{exercise}\label{transfer-learning:exr:4}
Explain what is wrong with fitting \texttt{PCA} on the features of all images, then splitting into
training and test sets, and why the reported test accuracy is optimistic.
\solutionlink{transfer-learning:exr:4}
\end{exercise}
\end{notechapter}
